\chapter{Достаточные условия равномерной сходимости тригонометрического ряда Фурье.}
%----------------------------------------------------------------------------------------
% Определение 22.12 --- кусочно-непрерывно дифференцируемая функция
%----------------------------------------------------------------------------------------

\begin{defn} \label{df22:12}
Функция $f$ называется \textit{кусочно-непрерывно дифференцируемой\rindex{функция!кусочно-непрерывно дифференцируемая}} на отрезке $[a;\, b]$, если отрезок $[a;\, b]$ можно разбить на конечное число отрезков, на каждом из которых функция $f$ (возможно, после изменения значений в концах) непрерывно дифференцируема (см. определение~6.5).
\end{defn}

\begin{notion}
В концах этих отрезков $x_i$, $i = 0, 1, \ldots, N$, функция $f$ не обязана быть непрерывной, но может иметь разрывы только первого рода; доопределение в концах происходит соответствующими предельными значениями $f(x_i + 0)$ и $f(x_i - 0)$.

Кусочно-непрерывно дифференцируемую на $[a;\, b]$ функцию $f$ можно определить как функцию такую, что $f$ и $f'$ имеют на $[a;\, b]$ не более конечного числа точек разрыва, причём только первого рода. Производная $f'$ такой функции кусочно-непрерывна на отрезке $[a;\, b]$ (см. определение~12.2); в точках $x_i$ производная $f'$ может не существовать.
\end{notion}

%----------------------------------------------------------------------------------------
% Определение 22.13 --- кусочно-гладкая функция
%----------------------------------------------------------------------------------------

\begin{defn} \label{df22:13}
Функция $f$ называется \textit{кусочно-гладкой\rindex{функция!кусочно-гладкая}} на отрезке $[a;\, b]$, если она непрерывна и кусочно-непрерывно дифференцируема на $[a;\, b]$.
\end{defn}

%----------------------------------------------------------------------------------------
% Теорема 22.6 --- равномерная сходимость и почленное дифференцирование ряда Фурье
%----------------------------------------------------------------------------------------

\begin{thm} \label{th22:6}
Пусть функция $f$ имеет период $2l$ и является кусочно-гладкой на отрезке $[-l;\, l]$, $l > 0$. Тогда:
\begin{enumerate}
    \item её ряд Фурье
    $$
    \frac{a_0}{2} + \sum_{n=1}^{\infty} \left(a_n \cos\frac{\pi n x}{l} + b_n \sin\frac{\pi n x}{l}\right)
    $$
    равномерно сходится к $f(x)$ на всей числовой прямой;
    \item производная $f'$ имеет ряд Фурье
    $$
    \sum_{n=1}^{\infty} \frac{\pi n}{l} \left(-a_n \sin\frac{\pi n x}{l} + b_n \cos\frac{\pi n x}{l}\right),
    $$
    который получается формальным дифференцированием ряда Фурье $f$.
\end{enumerate}
\end{thm}

\begin{proof}
Докажем сначала вторую часть теоремы. В условиях теоремы~\ref{th22:6} ряд Фурье $f$ сходится к $f(x)$ в каждой точке по следствию~1 из признака Липшица (в каждой точке $f$ либо дифференцируема, либо непрерывна и имеет конечные односторонние производные).

Далее, $f'$ кусочно-непрерывна на $[-l;\, l]$, следовательно, интегрируема по Риману; поэтому $f' \in L_R^2[-l;\, l]$. Пусть $f'$ имеет ряд Фурье
$$
\frac{\alpha_0}{2} + \sum_{n=1}^{\infty} \left(\alpha_n \cos\frac{\pi n x}{l} + \beta_n \sin\frac{\pi n x}{l}\right).
$$

Напомним, что для кусочно-непрерывных функций сохраняется формула Ньютона"--~Лейбница, где $f$ "--- обобщённая первообразная для $f'$ (теорема~12.22), поэтому
$$
\alpha_0 = \frac{1}{l} \int_{-l}^{l} f'(x)\,dx = \frac{1}{l}\bigl(f(l) - f(-l)\bigr) = 0,
$$
так как $f$ имеет период $2l$.

Далее, формула интегрирования по частям сохраняется, если функции $u$ и $v$ непрерывны на отрезке интегрирования, а $u'$ и $v'$ кусочно-непрерывны на нём (замечание после теоремы~12.24), поэтому при $n = 1, 2, \ldots$
\begin{align*}
\alpha_n &= \frac{1}{l} \int_{-l}^{l} f'(x) \cos\frac{\pi n x}{l}\,dx = \\
&= \frac{1}{l} \cdot f(x) \cos\frac{\pi n x}{l}\bigg|_{-l}^{l} - \int_{-l}^{l} f(x) \cdot \left(-\frac{\pi n}{l}\right) \sin\frac{\pi n x}{l}\,dx = \\
&= \frac{\pi n}{l} \cdot \frac{1}{l} \int_{-l}^{l} f(x) \sin\frac{\pi n x}{l}\,dx = \frac{\pi n}{l}\, b_n
\end{align*}
(обынтегрированный член равен нулю, так как $f(l) = f(-l)$); аналогично $\beta_n = -\dfrac{\pi n}{l}\, a_n$, $n = 1, 2, \ldots$ \textit{Вторая часть теоремы доказана.}

Далее, по неравенству Коши"--~Буняковского:
$$
\sum_{k=1}^{n} |a_k| = \frac{l}{\pi} \cdot \sum_{k=1}^{n} |\beta_k| \cdot \frac{1}{k} \le \frac{l}{\pi} \cdot \sqrt{\sum_{k=1}^{n} \beta_k^2} \cdot \sqrt{\sum_{k=1}^{n} \frac{1}{k^2}}.
$$
Но ряд $\sum_{k=1}^{\infty} \beta_k^2$ сходится, так как $f' \in L_R^2[-l;\, l]$; применено неравенство Бесселя "--- теорема ~\ref{th22:2}. Ряд $\sum_{k=1}^{\infty} \dfrac{1}{k^2}$ также сходится; частичные суммы этих рядов ограничены, и
$$
\sum_{k=1}^{n} |a_k| \le C, \quad n = 1, 2, \ldots
$$
Но если ряд с неотрицательными членами имеет ограниченные частичные суммы, то он сходится (теорема~15.5). Итак, ряд $\sum_{n=1}^{\infty} |a_n|$ сходится, аналогично сходится ряд $\sum_{n=1}^{\infty} |b_n|$.

Так как при всех $x$
$$
\left|a_n \cos\frac{\pi n x}{l} + b_n \sin\frac{\pi n x}{l}\right| \le |a_n| + |b_n|, \quad n = 1, 2, \ldots,
$$
то по признаку Вейерштрасса равномерной сходимости рядов (теорема~16.6) ряд Фурье $f$ сходится равномерно на всей числовой прямой. Как уже отмечалось, в каждой точке $x$ он сходится к значению $f(x)$.
\end{proof}

\begin{notion}[Замечание 1]
В условиях теоремы~\ref{th22:6} ряд Фурье $f'$ не обязан сходиться, так как $f'$ всего лишь кусочно-непрерывная на $[-l;\, l]$ функция.

Для сходимости ряда Фурье мало кусочной непрерывности или даже непрерывности функции; существуют непрерывные периодические функции, ряды Фурье которых расходятся в некоторых точках (построение примеров выходит за рамки нашего курса).

Для сходимости и равномерной сходимости рядов Фурье нужны некоторые дифференциальные свойства (условие Липшица, наличие конечных односторонних производных, кусочная гладкость и т.д.). Грубой ошибкой является вывод о сходимости в точке ряда Фурье на основании непрерывности функции в этой точке.
\end{notion}

\begin{notion}[Замечание 2]
Равномерная сходимость на всей числовой прямой ряда Фурье и возможность его почленного дифференцирования "--- разные свойства функции. Теорема~\ref{th22:6} даёт очень грубые условия, достаточные для выполнения обоих этих свойств, но это не означает равносильность этих свойств.
\end{notion}